\subsection{Every fixed nontrivial coefficient truncation has unbounded cotangent support}

\paragraph{A-S09-12 (Proved; all fixed truncations $S/I_N$, $N\geq2$).}
Put $R_N=\Gamma_{\mathbb Z}(T)/I_N$, where $I_N$ is the full
weight tail starting at $N$, and let $N\geq2$.  For every odd
prime $p>3(N-1)$, every $k\geq1$, and every cellular free
simplicial PD $S$-algebra resolution $Q\to R_N$,
\[
 \pi_{2k-1}\!\left(\Omega^1_{\mathrm{PD}}(Q/S)\right)^{p^k}
          \otimes_{\mathbb Z}\mathbb Z_{(p)}\ne0.
 \tag{S09.21}
\]
The free cells, the resolution, and the differential module here
are in weight-graded simplicial PD algebras, with special
differentials $d\gamma_m(x)=\gamma_{m-1}(x)\,dx$.
Thus changing the cellular resolution of a fixed $R_N$ cannot
give an eventual fixed homological bound in row one.  The
cutoff $N$ is fixed before choosing $p$; the assertion does
not apply to the prime-dependent choice $N=p^e$ of C09.

\begin{proof}
We construct one resolution and its comparisons explicitly.
Write
\[
 M_N=R_N^+=\bigoplus_{j=1}^{N-1}\mathbb Zt_j,
 \qquad t_j=T^{[j]},\qquad G=\Gamma^+_{\mathbb Z}.
\]
Let $\alpha:G(M_N)\to M_N$ be the actual PD-algebra
structure of $R_N$.  The cotriple
\[
 Q_n=\mathbb Z\oplus G^{n+1}(M_N)\longrightarrow R_N
 \tag{S09.22}
\]
has faces given by substitution and $\alpha$, degeneracies
given by the unit of $G$, and structural map from $S$ sending
$T$ to the fully unary iterate of $t_1$.  The additive extra
degeneracy contracts its augmentation kernel.  Its degeneracy
maps insert unary layers into the integral divided-power
monomial basis.  Removing all wholly unary layers gives the
unique nondegenerate ancestor.  As in A-S09-09, adjoining
these remaining generators along their prescribed boundaries
is a cellular description over $S$.  Here its degree-zero
cells also include $t_2,\ldots,t_{N-1}$.

Let
\[
 K_N=N\operatorname{Bar}_\bullet(\mathrm{id},G,M_N),
 \qquad
 C_N=N\bigl(R_N\otimes_Q\Omega^1_{\mathrm{PD}}(Q/S)\bigr).
 \tag{S09.23}
\]
The first bar face of $K_N$ is arity-one projection, the
interior faces are substitution, and its last face uses the
actual map $\alpha$.  In degree $n$, before normalization,
the special differential module is
$R_N\otimes G^n(M_N)$ modulo the constant copy $R_N\,dt_1$.
The first face is special differentiation with coefficients
evaluated in $R_N$; the other faces are the bar faces just
specified.  The removed constant copy has weights at most
$N$.  In every weight each normalized complex is finite free,
and a normalized tree of degree $n$ and weight $w$ has
$n\leq w-1$.

\emph{The bar filtration.}
Filter $K_N^w$ increasingly by the number of bottom letters,
counting their divided-power multiplicities.  A letter $t_j$
counts once and has weight $j$.  The two counts are distinct.
Every substitution face preserves the number of bottom
letters.  Applying $\alpha$ to a bottom monomial of length
$b>1$ either gives zero or a scalar multiple of one letter
$t_j$, and therefore strictly decreases that number.
The associated graded differential is consequently the bar
differential with the trivial algebra structure on $M_N$:
only unary bottom operations survive.  This is an equality
of the displayed integral associated graded complexes, not
an identification asserted only in a homotopy category.

Reduce this finite filtration modulo $p$.  Its first homology
page is computed by the same opened Fu--Nuiten comparison
used in A-S09-11, now giving each letter $t_j$ the positive
multiweight $(j,1)$.  The first coordinate is the original
weight and the second is bottom-letter count.  Each letter
has homological degree zero.  A Lyndon word of length $L$
has degree $1-L$ and total original weight $b$ satisfying
\[
                     L\leq b\leq(N-1)L.
 \tag{S09.24}
\]
An operation sequence of length $r$, with the additional
index $e\in\{0,1\}$ allowed by that theorem, has original
weight $(1+e)p^r b$ and bottom-letter count
$(1+e)p^r L$.  Its bar homological degree is the negative
of the Lie degree in Theorem C.  If $r=0$, that degree is
\[
                          n=(1+e)L-1.                \tag{S09.25}
\]
For $r\geq1$, all sequence entries are less than one.
The alternative family with entries greater than one is
excluded by the theorem's upper bound on the terminal
entry, since $1-L\leq0$.  An entry zero can occur only
first.  Thus the sum of the entries is at most $-(r-1)$,
and, when $e=0$,
\[
                          n\geq L+2r-2.              \tag{S09.26}
\]

Take $w=p^k$.  Its oddness excludes $e=1$.  If $r=k$,
then $b=1$, so the word is the single letter $t_1$.
Exactly as in A-S09-11, it contributes one class in each
of degrees $2k-1$ and $2k$, both in the maximal possible
letter filtration $p^k$.

For every other class put $j=k-r\geq1$.  By (S09.24),
\[
 L\geq\left\lceil\frac{p^j}{N-1}\right\rceil
       \geq2j+2.
 \tag{S09.27}
\]
For the last inequality, $p>3(N-1)$ gives
$p^j/(N-1)>3p^{j-1}\geq2j+1$.
If $r\geq1$, (S09.26) now gives $n\geq2k$.
If $r=0$, (S09.25) gives $n\geq2k+1$.
Hence the first page has no homology in degree less than
$2k-1$, and exactly one class in degree $2k-1$.
That class has maximal letter filtration.  No subsequent
differential can enter it from a higher filtration, and no
differential can leave it for degree $2k-2$, which is zero
on the first page and remains zero.  The finite filtration
therefore proves
\[
 H_{2k-1}(K_N^{p^k}\otimes\mathbb F_p)=\mathbb F_p,
 \qquad
 H_n(K_N^{p^k}\otimes\mathbb F_p)=0\quad(n<2k-1).
 \tag{S09.28}
\]

\emph{The coefficient comparison.}
For $w>N$, coefficient augmentation is a chain surjection
\[
 \rho:C_N^w\longrightarrow K_N^w,\qquad
       a\otimes x\longmapsto\epsilon(a)x.
 \tag{S09.29}
\]
The first-face verification is the one in A-S09-11:
after coefficient augmentation special differentiation
retains exactly the unary outer monomials.  Its kernel
has a finite descending filtration by coefficient weight.
The associated quotients, for $a=1,\ldots,N-1$, are
precisely $K_N^{w-a}$.  Additional differentiated
coefficients strictly increase the coefficient weight;
all other faces preserve it.  This proves the assertion
about the associated graded differential directly.

For $w=p^k$, each $m=w-a$, $1\leq a<N$, is prime to
$p$.  In the first page of the bottom-letter filtration
of $K_N^m\otimes\mathbb F_p$, an operation sequence
must therefore have length zero.  Equation (S09.25) and
the weight bound give
\[
 n+1=(1+e)L\geq\frac{m}{N-1}
       \geq\frac{p^k}{N-1}-1
       >3p^{k-1}-1\geq2k.
\]
Thus all its homology in degrees $n<2k$ is zero.
The finite coefficient filtration and its long exact
sequences imply the same vanishing for
$\ker(\rho)\otimes\mathbb F_p$.  Since $p^k>N$, the
relative quotient $R_N\,dt_1$ has no component in this
weight.  Equations (S09.28)--(S09.29) give
\[
 H_{2k-1}(C_N^{p^k}\otimes\mathbb F_p)=\mathbb F_p,
 \qquad H_{2k-2}(C_N^{p^k}\otimes\mathbb F_p)=0.
 \tag{S09.30}
\]
The finite-free integral universal-coefficient argument
of A-S09-11 proves that
$H_{2k-1}(C_N^{p^k})\otimes\mathbb Z_{(p)}$ is nonzero.

\emph{Transfer to every stated resolution.}
First, for (S09.22), the cellular differential-module
filtration of A-S09-09 proves that
$\Omega^1_{\mathrm{PD}}(Q/S)\to
R_N\otimes_Q\Omega^1_{\mathrm{PD}}(Q/S)$ is a weak
equivalence.  Its associated cell comparisons are
$Q\otimes\mathbb Z[\Delta^n/\partial\Delta^n]\to
R_N\otimes\mathbb Z[\Delta^n/\partial\Delta^n]$.
They are weak equivalences, and the degreewise split
filtration and exact filtered colimits preserve this property.

Let $Q'\to R_N$ be another cellular free resolution.
Both augmentations are surjective in every simplex and
weight: their restrictions to $S\to R_N$ already are.
Their kernels are acyclic simplicial abelian groups.
An elementary filling argument constructs algebra maps
$u:Q\to Q'$ and $v:Q'\to Q$ over $R_N$, and homotopies
$vu\simeq\mathrm{id}_Q$ and $uv\simeq\mathrm{id}_{Q'}$
over $R_N$.  Here are the needed details of that argument.
At each new free simplex generator the known faces give
a map from the boundary of its simplex into the target.
After choosing a lift of its prescribed augmentation,
the difference is a boundary problem in the augmentation
kernel.  Normalization splits off the degenerate terms;
the remaining obstruction is a cycle in degree $n-1$
of that acyclic kernel and is filled by a degree-$n$
element.  The chosen value of the free generator extends
uniquely as a PD-algebra map.  Repeating along the cells
constructs $u$ and $v$.  The same procedure on the
simplex prisms $\Delta^n\times\Delta^1$, relative to
their boundaries and two ends, constructs the two
homotopies; each prism has a finite simplex decomposition,
so it is the identical boundary-filling problem.

Apply the functor
$B\mapsto R_N\otimes_B\Omega^1_{\mathrm{PD}}(B/S)$
to these maps and homotopies.  Every component of a
simplicial algebra homotopy is a PD-algebra map over
$S$ and $R_N$; special differentiation therefore gives
the corresponding $R_N$-linear homotopy and preserves
its simplicial identities.  We obtain an actual simplicial
homotopy equivalence between the two base-changed
differential modules.  The same cellular filtration gives
the weak equivalence from the un-base-changed module
for $Q'$ as well.  These explicit comparisons transfer
the nonzero group just proved to $Q'$, establishing
(S09.21) in its full stated scope.
\end{proof}

The supplementary integer certificate
\texttt{A\_cosimplicial\_pd\_truncation.py} checks 7,209
individual face identities and 1,193 filtration identities
at $N=2,3,4$ in weights through six.  It also constructs
the entire weight-seven complex for $N=3$ through degree
three, whose dimensions are $(0,11,88,262)$, and checks
350 columns of $d^2=0$.  The differential ranks modulo
seven are $(0,10,77)$, giving the predicted one-dimensional
$H_1$ in this instance.  These checks verify the explicit
chain formulas; the all-$N$ assertion follows from the
filtration proof, not from extending a finite pattern.
